\section{Database / Data Exposure}
\label{sec:db}

\subsection{Finding IC-03 --- \fillme{e.g. "Cleartext Database Credentials Reachable via RCE"}}

\begin{center}
\begin{tabularx}{\linewidth}{@{}p{3.8cm} X p{3.4cm} X@{}}
\toprule
\rowcolor{TableHead}
\multicolumn{4}{l}{\textcolor{white}{\bfseries\sffamily Finding Metadata}} \\
\midrule
\textbf{ID} & IC-03 & \textbf{Severity} & \critical \\
\rowcolor{TableAlt}
\textbf{Category} & \fillme{Credential exposure / data at risk} & \textbf{Status} & \fillme{Reachable now} \\
\textbf{Backends} & \fillme{e.g. PostgreSQL, Redis} & \textbf{Action} & Rotate all credentials \\
\bottomrule
\end{tabularx}
\end{center}

\subsubsection{Evidence}
\begin{enumerate}
  \item \textbf{Credentials are stored in cleartext.} \fillme{Describe the config file / location
        and how many accounts/hosts are affected. Do not include the actual secret values.
        \emph{Values are withheld from this report per the rules of engagement.}}
  \item \textbf{Reachability.} \fillme{Explain how the credential file becomes readable given the
        other artifacts in this report (e.g. via the webshell in IC-01).}
  \item \textbf{Capability.} \fillme{State what an actor can do with OS access plus these
        credentials.}
\end{enumerate}

\begin{criticalbox}
\fillme{One paragraph: because a live foothold and reachable cleartext credentials coexist, every
credential reachable from this host must be treated as burned. State the rotation scope and the
follow-up review needed (unauthorised accounts, unexpected logons, schema changes).}
\end{criticalbox}

\subsubsection{Data at Risk}
\fillme{Name the system this supports and list, at a category level (not specific records), what
kinds of sensitive data are exposed if this path is exploited -- e.g. personal records, financial
data, credentials for other systems.}
\begin{itemize}
  \item \fillme{Data category 1}
  \item \fillme{Data category 2}
\end{itemize}
